% figures14.tex -- Chapter 14 (Space Geometry) figures redrawn in TikZ.
% Proportions follow the printed figures (iPad scan15 pp.3-19 = book pp.236-249);
% exact coordinates are reconstructions.
%
% Projection conventions copied from the book:
%   * a plane is a parallelogram, long axis horizontal, receding edge up-right;
%   * hidden edges are dashed (style "hid" below -- part of the GIVEN object,
%     so it stays black in the colour edition; it is not an "aux" line);
%   * solids are drawn in oblique projection with the depth axis up-right.

\tikzset{
  hid/.style  = {fig, dashed, dash pattern=on 2.6pt off 2.2pt},
  arw/.style  = {fig, ->, >=latex},
  % A dihedral/plane-angle arc, as the book strikes it.
  farc/.style = {fig, line width=0.6pt},
}
% Point dots are \dt{} (style/brumfiel.sty): a node of fixed 3.2pt diameter, so
% it prints the same size whatever `scale=' the picture carries.  The book's
% points are heavy round discs -- never arrowheads (Caleb, 2026-08-17).

% \XIVplane{name}{ox}{oy}{width}{dx}{dy} -- a plane as a parallelogram.
% Defines <name>sw, <name>se, <name>ne, <name>nw and strokes the outline.
\newcommand{\XIVplane}[6]{%
  \coordinate (#1sw) at (#2,#3);
  \coordinate (#1se) at (#2+#4,#3);
  \coordinate (#1ne) at (#2+#4+#5,#3+#6);
  \coordinate (#1nw) at (#2+#5,#3+#6);
  \draw[fig] (#1sw) -- (#1se) -- (#1ne) -- (#1nw) -- cycle;}

% \XIVplanec{name}{ox}{oy}{width}{dx}{dy} -- same, corners only, no stroke.
\newcommand{\XIVplanec}[6]{%
  \coordinate (#1sw) at (#2,#3);
  \coordinate (#1se) at (#2+#4,#3);
  \coordinate (#1ne) at (#2+#4+#5,#3+#6);
  \coordinate (#1nw) at (#2+#5,#3+#6);}

% ---------------------------------------------------------------- Fig 14-1
% Five familiar solids: two intersecting planes, cube, sphere, pentagonal
% prism, square pyramid.  Scan: the cube carries no hidden edges; the prism
% and the pyramid do.
\newcommand{\FIGXIVONE}{\begin{bkfigure}[\linewidth]
\begin{tikzpicture}[scale=0.760]
  % (1) two intersecting planes.  The book draws two long parallelogram
  %     "boards" that genuinely CROSS near their lower-right ends, leaving the
  %     two little spikes where each corner pokes out past the other.  (The old
  %     version built both on one common line, which put three near-coincident
  %     long edges down the middle -- the "smear" Caleb flagged.)
  \coordinate (oA) at (1.05,2.90);  \coordinate (dA) at (1.20,-2.60);
  \coordinate (wA) at (0.55,0.18);
  \coordinate (oB) at (0.05,1.75);  \coordinate (dB) at (2.55,-1.55);
  \coordinate (wB) at (0.42,0.62);
  \draw[fig] (oA) -- ($(oA)+(dA)$) -- ($(oA)+(dA)+(wA)$) -- ($(oA)+(wA)$)
             -- cycle;
  \draw[fig] (oB) -- ($(oB)+(dB)$) -- ($(oB)+(dB)+(wB)$) -- ($(oB)+(wB)$)
             -- cycle;
  % (2) cube
  \begin{scope}[xshift=3.55cm, yshift=0.45cm]
    \draw[fig] (0,0) -- (1.45,0) -- (1.45,1.45) -- (0,1.45) -- cycle;
    \draw[fig] (0,1.45) -- (0.62,1.95) -- (2.07,1.95) -- (1.45,1.45);
    \draw[fig] (1.45,0) -- (2.07,0.5) -- (2.07,1.95);
  \end{scope}
  % (3) sphere.  The book shades it with a DENSE fine dot screen filling a
  %     broad annulus just inside the rim, heaviest at the upper left and
  %     thinning round the bottom, leaving a clean highlight low and right.
  %     `pattern=dots' is far too sparse for that, so the screen is drawn.
  \begin{scope}[xshift=7.05cm, yshift=1.5cm]
    \begin{scope}[even odd rule]
      \clip (0,0) circle (0.86);
      \clip (0.86,-0.86) rectangle (-0.86,0.86) (0.26,-0.22) circle (0.72);
      \foreach \i in {0,...,31}{\foreach \j in {0,...,31}{
        \fill ({-0.87+0.056*\i+0.028*mod(\j,2)},{-0.87+0.056*\j})
              circle (0.30pt);}}
    \end{scope}
    \draw[fig] (0,0) circle (0.86);
  \end{scope}
  % (4) rectangular box, standing on end and tilted, drawn open at the top --
  %     the book's fourth solid.  e1 = width, e2 = depth, e3 = length.
  \begin{scope}[xshift=8.95cm, yshift=2.40cm]
    \coordinate (e1) at (1.02,0.09);
    \coordinate (e2) at (0.60,-0.83);
    \coordinate (e3) at (-1.02,-2.33);
    \coordinate (TL) at (0,0);
    \coordinate (TT) at ($(TL)+(e1)$);
    \coordinate (TR) at ($(TL)+(e1)+(e2)$);
    \coordinate (TB) at ($(TL)+(e2)$);
    \coordinate (BL) at ($(TL)+(e3)$);
    \coordinate (BT) at ($(TT)+(e3)$);
    \coordinate (BR) at ($(TR)+(e3)$);
    \coordinate (BB) at ($(TB)+(e3)$);
    \draw[fig] (TL) -- (TT) -- (TR) -- (TB) -- cycle;   % top face (rim)
    \draw[fig] (TL) -- (BL) -- (BB) -- (TB);            % near-left face
    \draw[fig] (BB) -- (BR) -- (TR);                    % near-right face
    \draw[hid] (BT) -- (TT);                            % hidden far edge
    \draw[hid] (BT) -- (BL);  \draw[hid] (BT) -- (BR);  % hidden bottom edges
    % the inner lip: the box reads as open at the top
    \coordinate (c) at ($(TL)!0.5!(TR)$);
    \draw[fig] ($(c)!0.80!(TL)$) -- ($(c)!0.80!(TB)$) -- ($(c)!0.80!(TR)$);
  \end{scope}
  % (5) square pyramid.  The book draws the apex-to-rear-vertex edge SOLID and
  %     fans the two hidden base edges from its foot as straight dashes.
  \begin{scope}[xshift=11.15cm, yshift=0.40cm]
    \coordinate (b1) at (0,0.34); \coordinate (b2) at (1.0,0);
    \coordinate (b3) at (2.0,0.34); \coordinate (b4) at (1.0,0.68);
    \coordinate (ap) at (1.0,1.92);
    \draw[fig] (b1) -- (b2) -- (b3);
    \draw[fig] (ap) -- (b1); \draw[fig] (ap) -- (b2); \draw[fig] (ap) -- (b3);
    \draw[fig] (ap) -- (b4);
    \draw[hid] (b3) -- (b4) -- (b1);
  \end{scope}
\end{tikzpicture}
\figcap{14--1}
\end{bkfigure}}

% ---------------------------------------------------------------- Fig 14-2
% Pictorial: a swept-wing jet above a sedan on a hatched ground line.  This is
% an illustration, not a geometric construction -- see ch14-UNCERTAIN.md.
\newcommand{\FIGXIVTWO}{\begin{bkfigure}[0.80\linewidth]
\begin{tikzpicture}[scale=0.806]
  % ---- jet, three-quarter view, nose to the right
  \begin{scope}[yshift=1.55cm]
    % fuselage
    \draw[fig] (0.30,0.42) -- (0.95,0.30) -- (4.55,0.28) -- (4.95,0.40)
               -- (4.95,0.56) -- (4.52,0.66) -- (1.05,0.70) -- (0.42,0.62)
               -- cycle;
    % canopy
    \draw[fig] (3.05,0.70) -- (3.35,0.94) -- (4.00,0.94) -- (4.22,0.70);
    \draw[fig] (3.42,0.70) -- (3.60,0.94);
    % swept main wing (far) and near wing
    \draw[fig] (1.85,0.70) -- (1.35,1.30) -- (2.55,1.24) -- (2.90,0.70);
    \draw[fig] (2.05,0.28) -- (1.50,-0.34) -- (2.75,-0.28) -- (3.05,0.28);
    % tailplane + fin
    \draw[fig] (0.62,0.70) -- (0.28,1.10) -- (1.05,1.06) -- (1.28,0.70);
    \draw[fig] (0.55,0.70) -- (0.45,1.42) -- (1.02,1.40) -- (1.30,0.70);
    % nose probe
    \draw[fig] (4.95,0.48) -- (5.35,0.48);
  \end{scope}
  % ---- sedan
  \begin{scope}[xshift=4.05cm]
    \draw[fig] (0.30,0.30) -- (0.36,0.62) -- (0.95,0.66) -- (1.55,1.02)
               -- (2.95,1.04) -- (3.30,0.68) -- (4.15,0.62) -- (4.22,0.30);
    \draw[fig] (0.30,0.30) -- (4.22,0.30);
    \draw[fig] (1.62,0.94) -- (2.22,0.96) -- (2.22,0.66);
    \draw[fig] (2.36,0.96) -- (2.90,0.96) -- (3.16,0.68);
    \draw[fig] (0.95,0.66) -- (3.30,0.68);
    \draw[fig] (1.28,0.30) circle (0.24);
    \draw[fig] (3.42,0.30) circle (0.24);
  \end{scope}
  % ---- ground line with hatching
  \draw[fig] (0.55,0.06) -- (9.20,0.06);
  \foreach \x in {0.75,1.05,...,9.05}
    {\draw[fig, line width=0.4pt] (\x,0.06) -- (\x-0.22,-0.16);}
\end{tikzpicture}
\figcap{14--2}
\end{bkfigure}}

% ---------------------------------------------------------------- Fig 14-3
\newcommand{\FIGXIVTHREE}{\begin{bkfigure}[0.52\linewidth]
\begin{tikzpicture}[scale=0.910]
  \XIVplane{P}{0}{0}{3.5}{0.95}{1.30}
  \dt{1.30,0.52}
  \dt{2.55,0.86}
  \dt{1.95,0.30}
\end{tikzpicture}
\figcap{14--3}
\end{bkfigure}}

% ------------------------------------------------------------ Fig 14-4, 14-5
\newcommand{\FIGXIVFOURFIVE}{\begin{bkfigure}[\linewidth]
\begin{minipage}{0.47\linewidth}\centering
\begin{tikzpicture}[scale=0.793]
  \XIVplane{P}{0}{0}{3.5}{0.95}{1.30}
  \dt{1.15,0.46}
  \dt{2.30,0.78}
  \dt{1.85,0.26}
  \dt{2.20,1.52}   % the fourth point, off the plane
\end{tikzpicture}
\figcap{14--4}
\end{minipage}\hfill
\begin{minipage}{0.47\linewidth}\centering
\begin{tikzpicture}[scale=0.793]
  \XIVplane{P}{0}{0}{3.5}{0.95}{1.30}
  \coordinate (Pp) at (2.05,0.55);
  \coordinate (Lu) at (3.10,2.00);           % upper end of l, above pi
  \coordinate (Lo) at ($(Pp)!-1.25!(Lu)$);   % lower end -- on the SAME line
  % The book hides the stretch of l that lies BEHIND the plane: solid until
  % l crosses the near edge, dashed from there to P, solid above P.
  \coordinate (Le) at (intersection of Lo--Lu and Psw--Pse);
  \draw[fig] (Lo) -- (Le);
  \draw[hid] (Le) -- (Pp);
  \draw[fig] (Pp) -- (Lu);
  \dt{Pp}
  \node[vlab, left=3pt] at (Pp) {$P$};
  \node[vlab, right=3pt] at (Lu) {$l$};
\end{tikzpicture}
\figcap{14--5}
\end{minipage}
\end{bkfigure}}

% ------------------------------------------------------------ Fig 14-6, 14-7
\newcommand{\FIGXIVSIXSEVEN}{\begin{bkfigure}[\linewidth]
\begin{minipage}{0.40\linewidth}\centering
\begin{tikzpicture}[scale=1.144]
  \XIVplane{P}{0}{0}{2.9}{0.85}{1.20}
  % PQ is lifted clear of the bottom edge so that "l" fits underneath it
  \coordinate (Pp) at (0.72,0.80); \coordinate (Qq) at (2.05,0.60);
  \draw[fig] (Pp) -- (Qq);
  \dt{Pp} \dt{Qq}
  \dt{2.26,0.96}
  \node[vlab, above left=3pt] at (Pp) {$P$};
  \node[vlab, right=3.5pt] at (Qq) {$Q$};
  \node[vlab, right=3.5pt] at (2.26,0.96) {$R$};
  \node[vlab, below=3pt] at ($(Pp)!0.5!(Qq)$) {$l$};
\end{tikzpicture}
\figcap{14--6}
\end{minipage}\hfill
\begin{minipage}{0.56\linewidth}\centering
\begin{tikzpicture}[scale=1.10]
  % Theorem 14-3: infinitely many planes contain one line.  The book draws the
  % axis l nearly end-on, so the four planes are parallelograms built on the
  % SAME short segment AB of l and fanned about it like a deck of cards; each
  % has real area, and its far corner is a heavy dot on the dashed transversal
  % m.  (The old drawing made the planes slivers and left l undrawn.)
  \coordinate (Av) at (-1.450,1.167);        % upper common vertex, on l
  \coordinate (Bv) at (0,0);                 % lower common vertex, on l
  \coordinate (d1) at (-2.917,1.450);        % first dot on m
  \coordinate (dm) at (2.867,0.650);         % direction of m
  \foreach \t in {0, 0.53, 0.695, 1.0}{
    \coordinate (D) at ($(d1)+\t*(dm)$);
    \coordinate (W) at ($(D)-(Av)$);
    \draw[fig] (Av) -- (D) -- ($(Bv)+(W)$) -- (Bv);}
  \draw[fig] (Av) -- (Bv) -- (0.233,-0.283);       % the axis l, overshooting B
  \draw[hid] ($(d1)-0.16*(dm)$) -- ($(d1)+1.30*(dm)$);   % m, dashed throughout
  \foreach \t in {0, 0.53, 0.695, 1.0}{\dt{$(d1)+\t*(dm)$}}
  \node[vlab, right=2pt] at ($(d1)+1.30*(dm)$) {$m$};
  \node[vlab, below right=1pt] at (0.233,-0.283) {$l$};
\end{tikzpicture}
\figcap{14--7}
\end{minipage}
\end{bkfigure}}

% ------------------------------------------------------------ Fig 14-8, 14-9
\newcommand{\FIGXIVEIGHTNINE}{\begin{bkfigure}[\linewidth]
\begin{minipage}{0.47\linewidth}\centering
\begin{tikzpicture}[scale=1.259]
  \XIVplane{P}{0}{0}{3.0}{0.75}{1.25}
  % l_1 runs Q -> E1, l_2 runs P -> E2; R is their true crossing.  The book
  % ends each segment on a HEAVY ROUND DOT at P and at Q -- no arrowheads --
  % and marks R with the largest dot of the three.
  \coordinate (Pp) at (1.10,1.00); \coordinate (Qq) at (0.95,0.45);
  \coordinate (E1) at (2.62,1.06); \coordinate (E2) at (2.68,0.36);
  \draw[fig] (Pp) -- (E2);
  \draw[fig] (Qq) -- (E1);
  \coordinate (R) at (intersection of Pp--E2 and Qq--E1);
  \dt{Pp}\dt{Qq}\dt{R}
  \node[vlab, left=4pt] at (Pp) {$P$};
  \node[vlab, left=4pt] at (Qq) {$Q$};
  \node[vlab, below=3pt] at (R) {$R$};
  \node[vlab, right=2.5pt] at (E1) {$l_1$};
  \node[vlab, right=2.5pt] at (E2) {$l_2$};
\end{tikzpicture}
\figcap{14--8}
\end{minipage}\hfill
\begin{minipage}{0.47\linewidth}\centering
\begin{tikzpicture}[scale=1.44]
  % Theorem 14-5.  The book draws an OPEN-BOOK hinge: the two planes share one
  % COMPLETE edge and swing away from each other, so nothing crosses anything.
  % (The old drawing had them straddle the line, so each plane ran out past the
  % other on both sides and the middle filled with a lens-shaped tangle.)
  \coordinate (I1) at (0,0);            % bottom-left end of the common edge
  \coordinate (I2) at (1.567,1.217);    % upper-right end of it
  \coordinate (ua) at (-0.333,1.017);   % pi_1 swings up-left
  \coordinate (ub) at (1.183,0.333);    % pi_2 swings out to the right
  \draw[fig] (I1) -- (I2) -- ($(I2)+(ua)$) -- ($(I1)+(ua)$) -- cycle;
  \draw[fig] (I1) -- (I2) -- ($(I2)+(ub)$) -- ($(I1)+(ub)$) -- cycle;
  \coordinate (Pp) at ($(I1)!0.31!(I2)$);
  \coordinate (Qq) at ($(I1)!0.81!(I2)$);
  \dt{Pp} \dt{Qq}
  \node[vlab] at ($(Pp)+(-0.28,0.24)$) {$P$};
  \node[vlab] at ($(Qq)+(-0.28,0.24)$) {$Q$};
  \node[vlab] at (0.10,0.95) {$\pi_1$};
  \node[vlab] at (1.25,0.68) {$\pi_2$};
\end{tikzpicture}
\figcap{14--9}
\end{minipage}
\end{bkfigure}}

% --------------------------------------------------------------- Fig 14-10
% Dihedral angle: the two half planes are solid, the rest of each plane dashed.
\newcommand{\FIGXIVTEN}{\begin{bkfigure}[0.60\linewidth]
\begin{tikzpicture}[scale=1.620]
  \coordinate (E1) at (0.55,0.35);  \coordinate (E2) at (2.95,1.35);  % edge l
  \coordinate (u) at (-0.55,0.85);   % offset of half plane 1 (up-left)
  \coordinate (v) at (0.75,-0.55);   % offset of half plane 2 (down-right)
  % half plane pi_1^+ (solid)
  \draw[fig] (E1) -- ($(E1)+(u)$) -- ($(E2)+(u)$) -- (E2);
  % half plane pi_2^+ (solid)
  \draw[fig] (E1) -- ($(E1)+(v)$) -- ($(E2)+(v)$) -- (E2);
  % the completing half planes, dashed
  \draw[hid] (E1) -- ($(E1)-(u)$) -- ($(E2)-(u)$) -- (E2);
  \draw[hid] (E1) -- ($(E1)-(v)$) -- ($(E2)-(v)$) -- (E2);
  \draw[fig] (E1) -- (E2);
  \draw[fig] (E2) -- (3.45,1.55);
  \draw[fig] (E1) -- (0.15,0.18);            % l runs past both vertices
  \node[vlab, above] at (3.45,1.55) {$l$};
  \node[vlab] at (1.25,1.13) {$\pi_1^+$};
  \node[vlab] at (1.90,0.48) {$\pi_2^+$};
\end{tikzpicture}
\figcap{14--10}
\end{bkfigure}}

% ---------------------------------------------------------------- Fig 14-11
% Split out of the old \FIGXIVELEVENTWELVE block: Ex.~14-3 \#1 cites 14-11 and
% \#2 cites 14-12, so each now stands on its own inside the exercise list.
\newcommand{\FIGXIVELEVEN}{\begin{bkfigure}[0.62\linewidth]
\begin{tikzpicture}[scale=1.55]
  % tetrahedron seen nearly edge-on: L apex at left, vertical edge T-B at right
  \coordinate (L) at (0,0.62); \coordinate (T) at (2.32,1.42);
  \coordinate (Bv) at (2.32,0.02); \coordinate (R) at (2.86,0.72);
  \draw[fig] (L) -- (T) -- (Bv) -- cycle;
  \draw[fig] (T) -- (R) -- (Bv);
  \draw[hid] (L) -- (R);
\end{tikzpicture}
\figcap{14--11}
\end{bkfigure}}

% ---------------------------------------------------------------- Fig 14-12
\newcommand{\FIGXIVTWELVE}{\begin{bkfigure}[0.46\linewidth]
\begin{tikzpicture}[scale=1.35]
  % irregular polyhedron (surface), no hidden edges in the book's drawing
  \draw[fig] (0.62,0.10) -- (0.06,1.02) -- (0.30,2.02) -- (1.28,2.52)
             -- (2.10,2.18) -- (2.30,1.08) -- (1.72,0.16) -- cycle;
  \draw[fig] (0.30,2.02) -- (1.42,1.42) -- (2.30,1.08);
  \draw[fig] (1.42,1.42) -- (1.28,2.52);
  \draw[fig] (1.42,1.42) -- (1.05,0.62) -- (1.72,0.16);
  \draw[fig] (1.05,0.62) -- (0.62,0.10);
\end{tikzpicture}
\figcap{14--12}
\end{bkfigure}}

% ----------------------------------------------------------- Fig 14-13, 14-14
\newcommand{\FIGXIVTHIRTEENFOURTEEN}{\begin{bkfigure}[\linewidth]
\begin{minipage}{0.52\linewidth}\centering
\begin{tikzpicture}[scale=1.05]
  % pi drawn as a rectangle (seen face on); angle AOB is a right angle
  \draw[fig] (0,0) rectangle (2.55,1.75);
  \coordinate (O) at (1.05,0.52);
  \coordinate (Aa) at (2.02,0.86);          % ray OA, direction (0.97,0.34)
  \coordinate (Bb) at (0.71,1.49);          % ray OB = OA turned exactly 90 deg
  \draw[arw] (O) -- (Aa); \draw[arw] (O) -- (Bb);
  \draw[hid] (O) -- ($(O)!-0.85!(Aa)$);      % the rest of the line of OA
  \draw[fig, line width=0.5pt]
        (1.3331,0.6192) -- (1.2339,0.9024) -- (0.9508,0.8031);
  \dt{O}
  \node[vlab] at (1.22,0.24) {$O$};
  \node[vlab] at (2.24,0.92) {$A$};
  \node[vlab] at (0.50,1.49) {$B$};
  \node[vlab] at (2.30,0.34) {$\pi$};
  % pi' drawn as a diamond
  \begin{scope}[xshift=2.62cm, yshift=-0.15cm]
    \draw[fig] (1.375,0) -- (2.75,1.375) -- (1.375,2.75) -- (0,1.375) -- cycle;
    \coordinate (Op) at (1.375,1.10);
    \coordinate (rp) at (0.50,1.10);        % r'
    \coordinate (rpp) at (1.375,2.21);      % r''  (perpendicular to r')
    % the book tips OA and OB with arrowheads but leaves r' and r'' plain
    \draw[fig] (Op) -- (rp); \draw[fig] (Op) -- (rpp);
    \draw[hid] (Op) -- ($(Op)!-1!(rp)$);     % the rest of the line of r'
    \draw[fig, line width=0.5pt] (1.03,1.10) -- (1.03,1.445) -- (1.375,1.445);
    \node[vlab] at (0.95,0.85) {$r'$};
    \node[vlab] at (1.68,1.72) {$r''$};
    \node[vlab] at (1.72,0.80) {$O'$};
    \node[vlab] at (1.375,0.42) {$\pi'$};
  \end{scope}
\end{tikzpicture}
\figcap{14--13}
\end{minipage}\hfill
\begin{minipage}{0.44\linewidth}\centering
\begin{tikzpicture}[scale=1.1]
  % Definition 14-3: l meets pi at P and is perpendicular to EVERY line of pi
  % through P, so the book draws two such lines, both passing through P.
  \XIVplane{P}{0}{0}{3.0}{1.15}{1.25}
  \coordinate (Pp) at (1.72,0.55);
  \coordinate (md) at (0.98,0.40);            % direction of m
  \coordinate (nd) at (-0.92,0.36);           % a second line of pi through P
  \draw[fig] (Pp) -- (1.72,2.05);
  \draw[fig] ($(Pp)-0.65*(md)$) -- ($(Pp)+1.15*(md)$);
  \draw[fig] ($(Pp)-0.85*(nd)$) -- ($(Pp)+0.80*(nd)$);
  % right-angle mark between l and m, sheared into the oblique projection
  \draw[fig, line width=0.5pt] ($(Pp)+(0,0.30)$)
        -- ($(Pp)+(0,0.30)+0.30*(md)$) -- ($(Pp)+0.30*(md)$);
  \dt{Pp}
  \node[vlab, above] at (1.72,2.05) {$l$};
  \node[vlab, below=1pt] at (Pp) {$P$};
  \node[vlab, right=2pt] at ($(Pp)+1.15*(md)$) {$m$};
  \node[vlab] at (0.72,0.26) {$\pi$};
\end{tikzpicture}
\figcap{14--14}
\end{minipage}
\end{bkfigure}}

% ----------------------------------------------------------- Fig 14-15, 14-16
\newcommand{\FIGXIVFIFTEENSIXTEEN}{\begin{bkfigure}[\linewidth]
\begin{minipage}{0.48\linewidth}\centering
\begin{tikzpicture}[scale=1.13]
  \XIVplane{P}{-0.35}{0}{3.35}{0.85}{1.20}
  % Theorem 14-6 / its hint: l_1 = PB and l_2 = PC are the two intersecting
  % lines of pi; B and C are chosen on them, Q is a point of the auxiliary
  % segment BC, and m is the line PQ.  Every one of l_1, l_2, m therefore runs
  % through P, and Q genuinely lies on BC.
  \coordinate (Pp) at (1.15,0.72);
  \coordinate (Aa) at (1.15,2.35);
  \coordinate (Ap) at (1.15,-0.91);          % A' -- P is the midpoint of AA'
  \coordinate (Bb) at (0.52,0.46);
  \coordinate (Cc) at (2.62,0.76);
  \coordinate (Qq) at ($(Bb)!0.52!(Cc)$);    % Q lies on BC
  \draw[fig] (1.15,2.72) -- (Aa);
  \draw[fig] (Aa) -- (Pp) -- (Ap);
  \draw[arw] (Pp) -- (Bb);                   % l_1
  \draw[arw] (Pp) -- (Cc);                   % l_2
  \draw[fig] (Pp) -- ($(Pp)!3.14!(Qq)$);     % m = PQ
  \draw[hid] (Bb) -- (Cc);
  \foreach \z in {Bb,Qq,Cc}{\draw[hid] (Aa) -- (\z); \draw[hid] (Ap) -- (\z);}
  \foreach \z in {Pp,Qq,Aa,Ap}{\dt{\z}}
  \node[vlab, above] at (1.15,2.72) {$l$};
  \node[vlab] at (1.42,2.42) {$A$};
  \node[vlab] at (1.46,-0.95) {$A'$};
  \node[vlab] at (0.90,0.89) {$P$};
  \node[vlab] at (1.74,0.39) {$Q$};
  \node[vlab] at (0.34,0.68) {$B$};
  \node[vlab] at (2.80,0.68) {$C$};
  \node[vlab] at (0.36,0.22) {$l_1$};
  \node[vlab] at (3.06,0.86) {$l_2$};
  \node[vlab] at (2.88,0.22) {$m$};
\end{tikzpicture}
\figcap{14--15}
\end{minipage}\hfill
\begin{minipage}{0.48\linewidth}\centering
\begin{tikzpicture}[scale=0.81]
  % (a) the plane perpendicular to l at P.  The book's plane is a parallelogram
  % with NEAR-VERTICAL sides; l runs solid to the heavy dot at P, is hidden
  % (dashed) for exactly the stretch that lies behind the plane -- from P to the
  % point where it leaves the parallelogram -- and solid again beyond.  The old
  % drawing put a wedge at P that read as an arrowhead and stranded the label.
  \coordinate (T1) at (0.30,2.05);  \coordinate (T2) at (1.35,1.50);
  \coordinate (T3) at (1.35,-0.10); \coordinate (T4) at (0.30,0.45);
  \draw[fig] (T1) -- (T2) -- (T3) -- (T4) -- cycle;
  \coordinate (L0) at (-0.55,1.32); \coordinate (L1) at (2.30,1.12);
  \coordinate (Pa) at ($(L0)!0.40!(L1)$);
  \coordinate (Xa) at (intersection of L0--L1 and T2--T3);
  \draw[fig] (L0) -- (Pa);
  \draw[hid] (Pa) -- (Xa);
  \draw[fig] (Xa) -- (L1);
  \dt{Pa}
  \node[vlab, left=2pt] at (L0) {$l$};
  \node[vlab, below=3pt] at (Pa) {$P$};
  \node[slab] at (0.90,-0.62) {(a)};
  % (b) l1, l2, l3 and l' all perpendicular to l at P.  The right-angle mark
  % belongs in the wedge between l and l' (it used to sit between l and l_3),
  % and every ray label sits just past the free end of its own ray.
  \begin{scope}[xshift=2.75cm]
    \coordinate (Pb) at (0.95,0.95);
    \coordinate (Lb) at (0.98,2.10);      % l
    \coordinate (Lp) at (1.88,1.88);      % l'
    \coordinate (L3) at (2.15,0.87);      % l_3
    \coordinate (L2) at (1.67,0.20);      % l_2
    \coordinate (L1b) at (0.10,0.26);     % l_1
    \draw[fig] (Pb) -- (Lb);
    \draw[fig] (Pb) -- (L3);
    \draw[fig] (Pb) -- (L1b);
    \draw[fig] (Pb) -- (L2);
    \draw[aux] (Pb) -- (Lp);
    \draw[farc] (0.9568,1.2099) -- (1.1406,1.3937) -- (1.1338,1.1338);
    \node[vlab, above] at (Lb) {$l$};
    \node[vlab, right=2pt] at (L3) {$l_3$};
    \node[vlab, below left=1pt] at (L1b) {$l_1$};
    \node[vlab, below right=1pt] at (L2) {$l_2$};
    \node[vlab, above right=0pt] at (Lp) {$l'$};
    \node[vlab] at (0.62,1.12) {$P$};
    \node[slab] at (1.05,-0.62) {(b)};
  \end{scope}
\end{tikzpicture}
\figcap{14--16}
\end{minipage}
\end{bkfigure}}

% ----------------------------------------------------------- Fig 14-17, 14-18
\newcommand{\FIGXIVSEVENTEENEIGHTEEN}{\begin{bkfigure}[\linewidth]
\begin{minipage}{0.47\linewidth}\centering
\begin{tikzpicture}[scale=1.55]
  % Theorem 14-8.  The book's plane stands close to upright (near-vertical left
  % and right edges, top and bottom edges sloping DOWN to the right) and the
  % hidden run of l goes from the piercing dot ALL THE WAY to the plane's right
  % edge -- ours used to go solid again while still inside the plane.
  \coordinate (T1) at (0.05,1.85);  \coordinate (T2) at (1.80,1.38);
  \coordinate (T3) at (1.74,-0.43); \coordinate (T4) at (-0.01,0.04);
  \draw[fig] (T1) -- (T2) -- (T3) -- (T4) -- cycle;
  \coordinate (L0) at (-0.64,0.98); \coordinate (L1) at (2.38,0.52);
  \coordinate (Pp) at ($(L0)!0.452!(L1)$);
  \coordinate (Xp) at (intersection of L0--L1 and T2--T3);
  \draw[fig] (L0) -- (Pp);
  \draw[hid] (Pp) -- (Xp);
  \draw[fig] (Xp) -- (L1);
  \dt{Pp}
  \dt{0.75,1.405}
  \node[vlab] at (0.95,1.25) {$Q$};
  \node[vlab, left=2pt] at (L0) {$l$};
\end{tikzpicture}
\figcap{14--17}
\end{minipage}\hfill
\begin{minipage}{0.47\linewidth}\centering
\begin{tikzpicture}[scale=1.180]
  \XIVplane{P}{0}{0}{2.6}{0.75}{1.05}
  \coordinate (F1) at (1.05,0.42); \coordinate (F2) at (1.95,0.66);
  \draw[fig] (F1) -- (1.05,1.72);
  \draw[fig] (F2) -- (1.95,1.96);
  \dt{F1} \dt{F2}
\end{tikzpicture}
\figcap{14--18}
\end{minipage}
\end{bkfigure}}

% --------------------------------------------------------------- Fig 14-19
% Left plane: two lines that are parallel, and two that intersect.
% Right pair: parallel planes, each carrying a line; the lines are not parallel.
\newcommand{\FIGXIVNINETEEN}{\begin{bkfigure}[0.90\linewidth]
\begin{tikzpicture}[scale=0.806]
  \XIVplane{A}{0}{0}{3.4}{0.95}{1.30}
  \draw[fig] (1.05,0.92) -- (2.55,1.10);
  \draw[fig] (1.15,0.62) -- (2.65,0.80);      % parallel to the line above
  \draw[fig] (1.72,0.10) -- (2.62,0.52);
  \draw[fig] (2.55,0.12) -- (1.85,0.52);
  % two parallel planes at the right
  \XIVplane{B}{4.55}{1.55}{3.4}{0.95}{1.30}
  \draw[fig] (6.30,2.05) -- (7.15,2.72);
  \XIVplane{C}{4.55}{-0.55}{3.4}{0.95}{1.30}
  \draw[fig] (5.65,0.12) -- (7.35,0.42);
\end{tikzpicture}
\figcap{14--19}
\end{bkfigure}}

% ---------------------------------------------------- Fig 14-20, 14-21, 14-22
% l1' and l2' (Fig 14-20) are drawn parallel to l1 and l2 (Fig 14-21).
\newcommand{\FIGXIVTWENTYTWENTYTWO}{\begin{bkfigure}[\linewidth]
\begin{minipage}{0.30\linewidth}\centering
\begin{tikzpicture}[scale=1.05]
  \coordinate (Pp) at (1.35,1.05);
  \draw[aux] ($(Pp)-(1.30,0.52)$) -- ($(Pp)+(1.30,0.52)$);   % l1' -- slope 0.40
  \draw[aux] ($(Pp)-(1.20,-0.72)$) -- ($(Pp)+(1.20,-0.72)$); % l2' -- slope -0.60
  \dt{Pp}
  \node[vlab] at (1.30,1.42) {$P$};
  \node[vlab] at (2.86,1.86) {$l_1'$};
  \node[vlab, right=4.5pt] at ($(Pp)+(1.20,-0.72)$) {$l_2'$};
\end{tikzpicture}
\figcap{14--20}
\end{minipage}\hfill
\begin{minipage}{0.34\linewidth}\centering
\begin{tikzpicture}[scale=0.88]
  \XIVplane{P}{0}{0}{3.2}{0.95}{1.45}
  \coordinate (Xx) at (1.85,0.75);
  % In the book both segments are SHORT and lie entirely inside the plane --
  % ours ran corner to corner and poked out through two edges.
  \draw[fig] ($(Xx)-(0.62,0.248)$) -- ($(Xx)+(0.62,0.248)$);   % l1, slope 0.40
  \draw[fig] ($(Xx)-(0.58,-0.348)$) -- ($(Xx)+(0.58,-0.348)$); % l2, slope -0.60
  \node[vlab] at (2.68,1.10) {$l_1$};
  \node[vlab] at (2.66,0.30) {$l_2$};
\end{tikzpicture}
\figcap{14--21}
\end{minipage}\hfill
\begin{minipage}{0.30\linewidth}\centering
\begin{tikzpicture}[scale=0.86]
  % Two parallel planes with their common perpendicular.  Each plane stands on
  % edge with VERTICAL front and back edges and top/bottom edges sloping down
  % to the right.  Dashing (the defect Caleb flagged, "fully inverted"): the
  % line is SOLID everywhere -- including the whole stretch BETWEEN the planes
  % -- and dashed only for the short run inside each plane, from its piercing
  % point back to that plane's far edge.  Each plane's near vertical edge is
  % broken with a small gap where the line passes in front of it; the far edge
  % is drawn straight through, since it stands in front of the hidden run.
  \coordinate (aL) at (0.10,2.60);  \coordinate (bL) at (0.93,2.125);
  \coordinate (cL) at (0.93,0.175); \coordinate (dL) at (0.10,0.65);
  \coordinate (aR) at (2.19,2.62);  \coordinate (bR) at (3.02,2.145);
  \coordinate (cR) at (3.02,0.195); \coordinate (dR) at (2.19,0.67);
  \draw[fig] (aL) -- (bL) -- (cL) -- (dL);
  \draw[fig] (aL) -- (0.10,1.205);  \draw[fig] (0.10,1.025) -- (dL);
  \draw[fig] (aR) -- (bR) -- (cR) -- (dR);
  \draw[fig] (aR) -- (2.19,1.268);  \draw[fig] (2.19,1.088) -- (dR);
  \draw[fig] (-0.395,1.1000) -- (0.465,1.1260);
  \draw[hid] (0.465,1.1260) -- (0.930,1.1400);
  \draw[fig] (0.930,1.1400) -- (2.576,1.1897);
  \draw[hid] (2.576,1.1897) -- (3.020,1.2031);
  \draw[fig] (3.020,1.2031) -- (3.330,1.2125);
\end{tikzpicture}
\figcap{14--22}
\end{minipage}
\end{bkfigure}}

% --------------------------------------------------------------- Fig 14-23
% Three parallel planes cut by l and l'; AB:BC = A'B':B'C' by construction.
\newcommand{\FIGXIVTWENTYTHREE}{\begin{bkfigure}[0.64\linewidth]
\begin{tikzpicture}[scale=1.60]
  \XIVplane{T}{0.55}{2.30}{2.85}{0.72}{0.62}
  \XIVplane{M}{0.28}{1.15}{2.85}{0.72}{0.62}
  \XIVplane{B}{0.00}{0.00}{2.85}{0.72}{0.62}
  % Two transversals.  In the book the LEFT-hand label at the top is l', but
  % its line runs down to the RIGHT through A', B', C'; the right-hand label is
  % l and its line runs down to the LEFT through A, B, C -- the two cross just
  % above the top plane.  Both are parameterised by height h (h = 1 in the top
  % plane, 0.5 in the middle, 0 in the bottom), so B and B' fall in the middle
  % plane and AB:BC = A'B':B'C' exactly, which is Theorem 14-12's claim.
  \coordinate (Aa) at (1.72,2.62);   \coordinate (Cc) at (0.60,0.32);
  \coordinate (Bb) at ($(Aa)!0.5!(Cc)$);
  \coordinate (Ap) at (2.05,2.62);   \coordinate (Cp) at (2.80,0.32);
  \coordinate (Bp) at ($(Ap)!0.5!(Cp)$);
  % Hidden-line runs, computed rather than guessed: the line is dashed exactly
  % where it passes BELOW a plane and its projection still falls inside that
  % plane's parallelogram.  That is a short break under each of A, B and C, and
  % solid everywhere else -- the reverse of the old drawing, which was dashed
  % almost end to end and so carried no depth information at all.
  \draw[fig] (2.101,3.402) -- (Aa);
  \draw[hid] (Aa) -- (1.564,2.300);
  \draw[fig] (1.564,2.300) -- (Bb);
  \draw[hid] (Bb) -- (1.004,1.150);
  \draw[fig] (1.004,1.150) -- (Cc);
  \draw[hid] (Cc) -- (0.444,0.000);
  \draw[fig] (0.444,0.000) -- (0.331,-0.232);
  \draw[fig] (1.795,3.402) -- (Ap);
  \draw[hid] (Ap) -- (2.154,2.300);
  \draw[fig] (2.154,2.300) -- (Bp);
  \draw[hid] (Bp) -- (2.529,1.150);
  \draw[fig] (2.529,1.150) -- (Cp);
  \draw[hid] (Cp) -- (2.837,0.205);
  \draw[fig] (2.837,0.205) -- (2.980,-0.232);
  \foreach \z in {Aa,Bb,Cc,Ap,Bp,Cp}{\dt{\z}}
  \node[vlab, left=3pt] at (Aa) {$A$};
  \node[vlab, left=3pt] at (Bb) {$B$};
  \node[vlab, below right=0pt] at (Cc) {$C$};
  \node[vlab, right=3pt] at (Ap) {$A'$};
  \node[vlab, right=3pt] at (Bp) {$B'$};
  \node[vlab] at (2.48,0.20) {$C'$};
  \node[vlab, above=1pt] at (1.795,3.402) {$l'$};
  \node[vlab, above=1pt] at (2.101,3.402) {$l$};
\end{tikzpicture}
\figcap{14--23}
\end{bkfigure}}

% ----------------------------------------------------------- Fig 14-24, 14-25
\newcommand{\FIGXIVTWENTYFOURTWENTYFIVE}{\begin{bkfigure}[\linewidth]
\begin{minipage}{0.50\linewidth}\centering
\begin{tikzpicture}[scale=1.34]
  \coordinate (V) at (1.72,3.05);
  \coordinate (Q1) at (0.42,0.72);   % P1
  \coordinate (Q2) at (0.95,1.28);   % P2
  \coordinate (Q3) at (2.62,1.22);   % P3
  \coordinate (Q4) at (1.42,0.42);   % P4
  \coordinate (Pt) at ($(Q4)!0.58!(Q3)$);   % P is a point OF the side P_4P_3
  % rays from V through each vertex, extended past the plane
  \foreach \z in {Q1,Q3,Q4,Pt}{\draw[fig] (V) -- ($(V)!1.42!(\z)$);}
  \draw[hid] (V) -- ($(V)!1.30!(Q2)$);
  % the polygon P1 P2 P3 P4
  \draw[fig] (Q1) -- (Q4) -- (Q3);
  \draw[hid] (Q3) -- (Q2) -- (Q1);
  \coordinate (Qq) at ($(V)!1.30!(Pt)$);
  \dt{Qq}
  \node[vlab, above=1.5pt] at (V) {$V$};
  \node[vlab, left=3.5pt] at (Q1) {$P_1$};
  \node[vlab] at (0.56,1.50) {$P_2$};
  \node[vlab, right=3.5pt] at (Q3) {$P_3$};
  \node[vlab, below left=3pt] at (Q4) {$P_4$};
  \node[vlab] at (2.44,0.84) {$P$};
  \node[vlab, right=3.5pt] at (Qq) {$Q$};
  \node[vlab] at (1.80,0.98) {$\pi$};
\end{tikzpicture}
\figcap{14--24}
\end{minipage}\hfill
\begin{minipage}{0.46\linewidth}\centering
\begin{tikzpicture}[scale=1.03]
  % Theorem 14-13's hint figure.  V sits at the upper LEFT and B far off to the
  % right, so VB and CB form the narrow wedge the book draws.  The rays V->A
  % and V->C overshoot A and C; the DASHED V->D stops dead at D.  (Ours used to
  % overshoot the dashed VD instead -- which read as a stray fragment -- and
  % carried no overshoot on VC at all.)
  \coordinate (Aa) at (0,0);
  \coordinate (V)  at (0.667,1.983);
  \coordinate (Bb) at (3.900,0.733);
  \coordinate (Cc) at (0.733,0.817);
  \coordinate (Dd) at ($(Aa)!0.402!(Bb)$);   % D lies on AB
  \draw[fig] (V) -- ($(V)!1.26!(Aa)$);
  \draw[fig] (V) -- ($(V)!2.17!(Cc)$);
  \draw[fig] (V) -- (Bb);
  \draw[hid] (V) -- (Dd);
  \draw[fig] (Aa) -- (Bb);
  \draw[fig] (Aa) -- (Cc); \draw[fig] (Cc) -- (Bb);
  \foreach \z in {V,Aa,Cc,Dd}{\dt{\z}}
  \node[vlab, above=2pt] at (V) {$V$};
  \node[vlab, left=3pt] at (Aa) {$A$};
  \node[vlab, right=3.5pt] at (Bb) {$B$};
  \node[vlab] at (0.54,0.99) {$C$};
  \node[vlab, below=3pt] at (Dd) {$D$};
\end{tikzpicture}
\figcap{14--25}
\end{minipage}
\end{bkfigure}}

% ----------------------------------------------------------- Fig 14-26, 14-27
\newcommand{\FIGXIVTWENTYSIXTWENTYSEVEN}{\begin{bkfigure}[\linewidth]
\begin{minipage}{0.44\linewidth}\centering
\begin{tikzpicture}[scale=1.407]
  % a tetrahedral angle flattened out: four triangles round C, which sits ON
  % the boundary -- the faces do not encircle the point.
  \coordinate (C)  at (1.52,1.30);
  \coordinate (T1) at (0.92,2.52);
  \coordinate (T2) at (2.38,2.52);
  \coordinate (Lf) at (0.15,1.30);
  \coordinate (B1) at (0.90,0.06);
  \coordinate (B2) at (2.32,0.05);
  \draw[fig] (T1) -- (T2) -- (C) -- (B2) -- (B1) -- (Lf) -- cycle;
  \draw[fig] (C) -- (T1); \draw[fig] (C) -- (Lf); \draw[fig] (C) -- (B1);
\end{tikzpicture}
\figcap{14--26}
\end{minipage}\hfill
\begin{minipage}{0.52\linewidth}\centering
\begin{tikzpicture}[scale=1.320]
  \XIVplane{P}{0}{0}{3.05}{1.15}{1.70}
  \coordinate (Q1) at (1.52,1.02);   % P1
  \coordinate (Q2) at (2.42,1.02);   % P2
  \coordinate (Q3) at (2.30,0.48);   % P3
  \coordinate (Q4) at (1.62,0.45);   % P4
  \coordinate (Q5) at (1.18,0.72);   % P5
  \fill[pattern=north east lines] (Q1) -- (Q2) -- (Q3) -- (Q4) -- (Q5) -- cycle;
  \draw[fig] (Q1) -- (Q2) -- (Q3) -- (Q4) -- (Q5) -- cycle;
  % labels sit outside the cell, on the outward normal from its centroid
  \node[vlab] at (1.285,1.265) {$P_1$};
  \node[vlab] at (2.735,1.170) {$P_2$};
  \node[vlab] at (2.610,0.315) {$P_3$};
  \node[vlab] at (1.435,0.255) {$P_4$};
  \node[vlab] at (0.830,0.710) {$P_5$};
\end{tikzpicture}
\figcap{14--27}
\end{minipage}
\end{bkfigure}}

% --------------------------------------------------------------- Fig 14-28
% The five regular polyhedra, laid out as in the book: tetrahedron,
% octahedron, icosahedron on the top row; cube and dodecahedron below.
\newcommand{\FIGXIVTWENTYEIGHT}{\begin{bkfigure}[\linewidth]
\begin{tikzpicture}[scale=1.076]
  % ---- tetrahedron
  \begin{scope}[xshift=0cm, yshift=3.30cm]
    \coordinate (t1) at (0.42,2.15); \coordinate (t2) at (0.00,0.42);
    \coordinate (t3) at (1.92,0.72); \coordinate (t4) at (1.05,0.00);
    \draw[fig] (t1) -- (t2) -- (t4) -- cycle;
    \draw[fig] (t1) -- (t3) -- (t4);
    \draw[hid] (t2) -- (t3);
    \node[slab, below] at (0.95,-0.34) {Tetrahedron};
  \end{scope}
  % ---- octahedron
  \begin{scope}[xshift=3.55cm, yshift=3.20cm]
    \coordinate (o1) at (1.15,2.42); \coordinate (o2) at (1.15,-0.10);
    \coordinate (o3) at (0.00,1.32); \coordinate (o4) at (2.32,1.28);
    \coordinate (o5) at (0.78,0.72); \coordinate (o6) at (1.62,1.72);
    \draw[fig] (o3) -- (o1) -- (o4) -- (o2) -- cycle;
    \draw[fig] (o3) -- (o5) -- (o4); \draw[fig] (o1) -- (o5) -- (o2);
    \draw[hid] (o3) -- (o6) -- (o4); \draw[hid] (o1) -- (o6) -- (o2);
    \node[slab, below] at (1.15,-0.52) {Octahedron};
  \end{scope}
  % ---- icosahedron (drawn as the book does: hexagonal silhouette)
  \begin{scope}[xshift=7.35cm, yshift=3.20cm]
    \foreach \i in {0,1,2,3,4,5}
      {\coordinate (h\i) at ({1.28+1.28*cos(30+60*\i)},{1.20+1.28*sin(30+60*\i)});}
    \draw[fig] (h0) -- (h1) -- (h2) -- (h3) -- (h4) -- (h5) -- cycle;
    % The book's projection: a horizontal chord across the top pair of
    % vertices carrying two interior nodes, a second chord across the bottom
    % pair carrying one, and no hidden edges drawn.
    \coordinate (u1) at ($(h2)!0.21!(h0)$);
    \coordinate (u2) at ($(h2)!0.81!(h0)$);
    \coordinate (ic) at ($(h3)!0.5!(h5)$);
    \draw[fig] (h2) -- (h0);
    \draw[fig] (h3) -- (h5);
    \draw[fig] (h1) -- (u1); \draw[fig] (h1) -- (u2);
    \draw[fig] (u1) -- (ic); \draw[fig] (u2) -- (ic);
    \draw[fig] (u1) -- (h3); \draw[fig] (u2) -- (h5);
    \draw[fig] (ic) -- (h4);
    \node[slab, below] at (1.28,-0.40) {Icosahedron};
  \end{scope}
  % ---- cube
  \begin{scope}[xshift=1.55cm, yshift=-0.20cm]
    \draw[fig] (0,0) -- (1.75,0) -- (1.75,1.75) -- (0,1.75) -- cycle;
    \draw[fig] (0,1.75) -- (0.68,2.35) -- (2.43,2.35) -- (1.75,1.75);
    \draw[fig] (1.75,0) -- (2.43,0.60) -- (2.43,2.35);
    \node[slab, below] at (1.20,-0.30) {Cube};
  \end{scope}
  % ---- dodecahedron
  \begin{scope}[xshift=5.65cm, yshift=-0.55cm]
    \foreach \i in {0,...,9}
      {\coordinate (d\i) at ({1.42+1.42*cos(18+36*\i)},{1.42+1.42*sin(18+36*\i)});}
    \draw[fig] (d0) -- (d1) -- (d2) -- (d3) -- (d4) -- (d5) -- (d6) -- (d7)
               -- (d8) -- (d9) -- cycle;
    \foreach \i in {0,...,4}
      {\coordinate (c\i) at ({1.42+0.62*cos(90+72*\i)},{1.42+0.62*sin(90+72*\i)});}
    \draw[fig] (c0) -- (c1) -- (c2) -- (c3) -- (c4) -- cycle;
    \draw[fig] (c0) -- (d2); \draw[fig] (c1) -- (d4); \draw[fig] (c2) -- (d6);
    \draw[fig] (c3) -- (d8); \draw[fig] (c4) -- (d0);
    \node[slab, below] at (1.42,-0.34) {Dodecahedron};
  \end{scope}
\end{tikzpicture}
\figcap{14--28}
\end{bkfigure}}

% ----------------------------------------------------------- Fig 14-29, 14-30
\newcommand{\FIGXIVTWENTYNINETHIRTY}{\begin{bkfigure}[\linewidth]
\begin{minipage}{0.50\linewidth}\centering
\begin{tikzpicture}[scale=1.10]
  % Theorem 14-17: a dihedral angle cut by two parallel planes.  The EDGE of
  % the dihedral is the left edge BL-TL of the "wall"; the near face F1 is the
  % solid quadrilateral, and the far face F2 lies behind it, so only its outer
  % edge shows -- dashed.  Each plane cuts the dihedral in a plane angle: at
  % the top it is the thin wedge TL-TR / TL-U, struck with an ARC and a
  % perpendicular tick; at the bottom the same wedge, its second side hidden
  % behind the wall and therefore dashed.  (The old drawing had two arrowheads
  % where those two arcs belong and no ticks at all.)
  \XIVplane{PB}{-1.30}{-0.85}{3.30}{0.95}{1.25}
  \XIVplane{PT}{-0.75}{1.55}{3.30}{0.95}{1.25}
  \coordinate (BLw) at (0,0);        \coordinate (TLw) at (0.55,2.40);
  \coordinate (BRw) at (1.25,0.12);  \coordinate (TRw) at (1.80,2.52);
  \coordinate (Ub)  at (1.20,0.34);  \coordinate (Ut)  at (1.75,2.74);
  \draw[fig] (TLw) -- (TRw) -- (BRw) -- (BLw) -- cycle;   % near face F1
  \draw[fig] (TLw) -- (Ut);                               % F2 trace, top plane
  \draw[hid] (Ut) -- (Ub);                                % F2 far edge, hidden
  \draw[hid] (BLw) -- (Ub);                       % F2 trace, bottom, hidden
  \draw[farc] ([shift=(5.48:0.78)]TLw) arc (5.48:15.82:0.78);
  \draw[farc] ([shift=(5.48:0.42)]BLw) arc (5.48:15.82:0.42);
  \coordinate (Qt) at ($(TLw)!0.66!(Ut)$);
  \draw[farc] (Qt) -- ($(TLw)!(Qt)!(TRw)$);
  \coordinate (Qb) at ($(BLw)!0.60!(Ub)$);
  \draw[farc] (Qb) -- ($(BLw)!(Qb)!(BRw)$);
\end{tikzpicture}
\figcap{14--29}
\end{minipage}\hfill
\begin{minipage}{0.46\linewidth}\centering
\begin{tikzpicture}[scale=1.05]
  % Definition 14-8.  A dihedral edge E0->E1 lies in the horizontal plane; one
  % face rises from it as the tilted flap, and the plane PERPENDICULAR to the
  % edge is the clean rectangle standing in front.  The edge passes BEHIND the
  % rectangle, so it is dashed, and it RISES to the right (ours used to run
  % flat along the rectangle's foot, which killed the depth reading); a
  % perpendicular tick marks the right angle at the rectangle's foot.
  \draw[fig] (-0.950,-0.523) -- (2.483,-0.110) -- (3.633,1.200)
             -- (0.183,0.800) -- cycle;                 % the horizontal plane
  \coordinate (E0) at (0,0);          \coordinate (E1) at (1.933,0.850);
  \coordinate (r)  at (1.517,0.150);  \coordinate (u)  at (0.067,1.140);
  \draw[fig] ($(E0)+(u)$) -- ($(E1)+(u)$) -- (E1);       % the flap
  \draw[hid] (E1) -- (E0);                              % the dihedral edge
  \draw[fig] (E0) -- ($(E0)+(r)$) -- ($(E0)+(r)+(u)$) -- ($(E0)+(u)$)
             -- cycle;                                  % the perpendicular
  \coordinate (Qe) at ($(E0)!0.25!(E1)$);
  \draw[farc] (Qe) -- ($(E0)!(Qe)!($(E0)+(r)$)$);
\end{tikzpicture}
\figcap{14--30}
\end{minipage}
\end{bkfigure}}

% ---------------------------------------------------------------- Fig 14-31
% Split out of the old \FIGXIVTHIRTYONETHIRTYTHREE block: Ex.~14-12 cites
% 14-31 in \#1, 14-32 in \#2 and 14-33 in \#3, so each stands on its own now.
\newcommand{\FIGXIVTHIRTYONE}{\begin{bkfigure}[0.74\linewidth]
\begin{tikzpicture}[scale=1.45]
  % A right hexagonal prism with its bases IN two parallel planes.  The old
  % hexagon put its vertices at 30+60i degrees, so the pairs (30,330) and
  % (150,210) shared an x -- their lateral edges landed on top of one another
  % and the solid read as a flat bar with a single dashed line down it.  These
  % vertex angles are the book's own irregular ones: no two share an x, so the
  % near side shows FOUR solid lateral edges and the far side TWO dashed.
  \XIVplane{PB}{0}{0}{4.00}{1.10}{1.05}
  \XIVplane{PT}{0}{2.60}{4.00}{1.10}{1.05}
  \foreach \i/\a in {0/-10, 1/36, 2/118, 3/181, 4/213, 5/301}
    {\coordinate (b\i) at ({2.20+1.00*cos(\a)},{0.52+0.44*sin(\a)});
     \coordinate (t\i) at ({2.20+1.00*cos(\a)},{3.12+0.44*sin(\a)});}
  % top base: seen from above, so the whole outline is visible
  \draw[fig] (t0) -- (t1) -- (t2) -- (t3) -- (t4) -- (t5) -- cycle;
  % bottom base: near half solid, far half hidden
  \draw[fig] (b3) -- (b4) -- (b5) -- (b0);
  \draw[hid] (b0) -- (b1) -- (b2) -- (b3);
  \foreach \i in {0,3,4,5}{\draw[fig] (b\i) -- (t\i);}
  \foreach \i in {1,2}{\draw[hid] (b\i) -- (t\i);}
\end{tikzpicture}
\figcap{14--31}
\end{bkfigure}}

% ---------------------------------------------------------------- Fig 14-32
\newcommand{\FIGXIVTHIRTYTWO}{\begin{bkfigure}[0.44\linewidth]
\begin{tikzpicture}[scale=1.10]
  % pentagonal pyramid.  q0 and q2 are the silhouette vertices and q1 the front
  % one, so those three lateral edges are solid; q3 and q4 lie behind, so their
  % lateral edges AND the three base edges joining them are dashed.  (The old
  % version drew the back edge q2--q3 and the lateral to q3 solid.)
  \coordinate (ap) at (1.42,2.72);
  \foreach \i in {0,...,4}
    {\coordinate (q\i) at ({1.42+1.32*cos(198+72*\i)},{0.62+0.46*sin(198+72*\i)});}
  \draw[fig] (q0) -- (q1) -- (q2);
  \draw[hid] (q2) -- (q3) -- (q4) -- (q0);
  \foreach \i in {0,1,2}{\draw[fig] (ap) -- (q\i);}
  \foreach \i in {3,4}{\draw[hid] (ap) -- (q\i);}
\end{tikzpicture}
\figcap{14--32}
\end{bkfigure}}

% ---------------------------------------------------------------- Fig 14-33
\newcommand{\FIGXIVTHIRTYTHREE}{\begin{bkfigure}[0.54\linewidth]
\begin{tikzpicture}[scale=1.15]
  % parallelepiped
  \coordinate (ea) at (2.65,0.32);   % length
  \coordinate (eb) at (1.15,0.95);   % depth
  \coordinate (ec) at (-0.62,1.05);  % height
  \coordinate (A0) at (0.62,0);
  \coordinate (A1) at ($(A0)+(ea)$);
  \coordinate (A2) at ($(A1)+(eb)$);
  \coordinate (A3) at ($(A0)+(eb)$);
  \coordinate (B0) at ($(A0)+(ec)$);
  \coordinate (B1) at ($(A1)+(ec)$);
  \coordinate (B2) at ($(A2)+(ec)$);
  \coordinate (B3) at ($(A3)+(ec)$);
  \draw[fig] (A0) -- (A1) -- (A2); \draw[fig] (A0) -- (B0);
  \draw[fig] (A1) -- (B1); \draw[fig] (A2) -- (B2);
  \draw[fig] (B0) -- (B1) -- (B2) -- (B3) -- cycle;
  \draw[hid] (A0) -- (A3) -- (A2);
  \draw[hid] (A3) -- (B3);
\end{tikzpicture}
\figcap{14--33}
\end{bkfigure}}

% ----------------------------------------------------------- Fig 14-34, 14-35
\newcommand{\FIGXIVTHIRTYFOURTHIRTYFIVE}{\begin{bkfigure}[\linewidth]
\begin{minipage}{0.48\linewidth}\centering
\begin{tikzpicture}[scale=1.046]
  % rectangular box, edges labelled l, w, h
  \coordinate (a0) at (0,0);       \coordinate (a1) at (2.85,0);
  \coordinate (a2) at (2.85,1.55); \coordinate (a3) at (0,1.55);
  \coordinate (d)  at (0.85,0.72);
  \draw[fig] (a0) -- (a1) -- (a2) -- (a3) -- cycle;
  \draw[fig] (a3) -- ($(a3)+(d)$) -- ($(a2)+(d)$) -- (a2);
  \draw[fig] (a1) -- ($(a1)+(d)$) -- ($(a2)+(d)$);
  \node[vlab] at (3.50,1.72) {$w$};
  \node[vlab] at (1.425,1.28) {$l$};
  \node[vlab] at (3.12,0.775) {$h$};
\end{tikzpicture}
\figcap{14--34}
\end{minipage}\hfill
\begin{minipage}{0.48\linewidth}\centering
\begin{tikzpicture}[scale=1.25]
  % Theorem 14-19: one polyhedron as the union of two, sharing the HATCHED
  % face.  The book's shared section is a big triangle spanning the whole near
  % face (ours was a small patch in one corner), its boundary heavy and solid,
  % its interior filled with dense steep hatching.  The one hidden vertex is F;
  % only its two hidden edges FA and FC are dashed.
  \coordinate (Av) at (-1.117,1.917);
  \coordinate (Bv) at ( 1.350,2.033);
  \coordinate (Cv) at ( 2.883,1.550);
  \coordinate (Dv) at ( 2.483,0.167);
  \coordinate (Ev) at ( 0,0);
  \coordinate (Fv) at (-0.017,1.167);
  \begin{scope}
    \clip (Bv) -- (Fv) -- (Ev) -- cycle;
    \foreach \k in {0,...,27}
      {\draw[fig, line width=0.35pt]
         ({-1.40+0.107*\k},-0.30) -- ({-1.40+0.107*\k+1.05},2.55);}
  \end{scope}
  \draw[fig] (Av) -- (Bv) -- (Cv) -- (Dv) -- (Ev) -- cycle;
  \draw[fig] (Bv) -- (Ev);  \draw[fig] (Bv) -- (Dv);
  \draw[fig] (Bv) -- (Fv) -- (Ev);
  \draw[hid] (Fv) -- (Av);  \draw[hid] (Fv) -- (Cv);
\end{tikzpicture}
\figcap{14--35}
\end{minipage}
\end{bkfigure}}

% ----------------------------------------------------------- Fig 14-36, 14-37
\newcommand{\FIGXIVTHIRTYSIXTHIRTYSEVEN}{\begin{bkfigure}[\linewidth]
\begin{minipage}{0.46\linewidth}\centering
\begin{tikzpicture}[scale=0.672]
  \coordinate (ea) at (2.85,0);      % length
  \coordinate (eb) at (0.95,0.62);   % depth
  \coordinate (ec) at (-0.55,0.95);  % height (leaning left)
  \coordinate (A0) at (0.62,0.10);
  \coordinate (A1) at ($(A0)+(ea)$);  \coordinate (A2) at ($(A1)+(eb)$);
  \coordinate (A3) at ($(A0)+(eb)$);
  \coordinate (B0) at ($(A0)+(ec)$);  \coordinate (B1) at ($(A1)+(ec)$);
  \coordinate (B2) at ($(A2)+(ec)$);  \coordinate (B3) at ($(A3)+(ec)$);
  \draw[fig] (A0) -- (A1) -- (A2); \draw[fig] (A0) -- (B0);
  \draw[fig] (A1) -- (B1); \draw[fig] (A2) -- (B2);
  \draw[fig] (B0) -- (B1) -- (B2) -- (B3) -- cycle;
  \draw[hid] (A0) -- (A3) -- (A2);  \draw[hid] (A3) -- (B3);
\end{tikzpicture}
\figcap{14--36}
\end{minipage}\hfill
\begin{minipage}{0.50\linewidth}\centering
\begin{tikzpicture}[scale=0.672]
  \coordinate (ea) at (3.15,0);
  \coordinate (eb) at (1.05,0.68);
  \coordinate (ec) at (-0.62,1.05);
  \coordinate (A0) at (0.68,0.10);
  \coordinate (A1) at ($(A0)+(ea)$);  \coordinate (A2) at ($(A1)+(eb)$);
  \coordinate (A3) at ($(A0)+(eb)$);
  \coordinate (B0) at ($(A0)+(ec)$);  \coordinate (B1) at ($(A1)+(ec)$);
  \coordinate (B2) at ($(A2)+(ec)$);  \coordinate (B3) at ($(A3)+(ec)$);
  \draw[fig] (A0) -- (A1) -- (A2); \draw[fig] (A0) -- (B0);
  \draw[fig] (A1) -- (B1); \draw[fig] (A2) -- (B2);
  \draw[fig] (B0) -- (B1) -- (B2) -- (B3) -- cycle;
  \draw[hid] (A0) -- (A3) -- (A2);  \draw[hid] (A3) -- (B3);
  % Altitude: a TRUE vertical between the two horizontal faces, as in the book
  % -- not a copy of the (leaning) lateral edge.  Its top sits on the edge
  % B1B2 of the top face; dropping by the height of ec lands it on the
  % hidden base face, so the segment is dashed.
  \coordinate (Ht) at ($(B1)!0.5!(B2)$);
  \coordinate (Hb) at ($(Ht)-(0,1.05)$);
  \draw[hid] (Ht) -- (Hb);
  \draw[fig, line width=0.4pt] ($(Ht)-(0.14,0)$) -- ($(Ht)+(0.14,0)$);
  \draw[fig, line width=0.4pt] ($(Hb)-(0.14,0)$) -- ($(Hb)+(0.14,0)$);
  \node[slab, right] at (4.95,1.72) {Altitude $h$};
  \node[slab, right] at (4.95,0.42) {Area $B$};
  \draw[arw] (4.90,1.72) -- ($(Ht)-(0.04,0.14)$);
  \draw[arw] (4.90,0.42) -- ($(A1)!0.55!(A2)+(0.12,-0.04)$);
\end{tikzpicture}
\figcap{14--37}
\end{minipage}
\end{bkfigure}}

% --------------------------------------------------------------- Fig 14-38
\newcommand{\FIGXIVTHIRTYEIGHT}{\begin{bkfigure}[0.94\linewidth]
\begin{tikzpicture}[scale=0.950]
  % (a) a parallelepiped of which a triangular prism is half
  \begin{scope}
    \coordinate (ea) at (2.55,0); \coordinate (eb) at (1.02,0.72);
    \coordinate (ec) at (0,1.42);
    \coordinate (A0) at (0,0);   \coordinate (A1) at ($(A0)+(ea)$);
    \coordinate (A2) at ($(A1)+(eb)$); \coordinate (A3) at ($(A0)+(eb)$);
    \coordinate (B0) at ($(A0)+(ec)$); \coordinate (B1) at ($(A1)+(ec)$);
    \coordinate (B2) at ($(A2)+(ec)$); \coordinate (B3) at ($(A3)+(ec)$);
    \draw[fig] (A0) -- (A1) -- (A2); \draw[fig] (A0) -- (B0);
    \draw[fig] (A1) -- (B1); \draw[fig] (A2) -- (B2);
    \draw[fig] (B0) -- (B1) -- (B2) -- (B3) -- cycle;
    \draw[hid] (A0) -- (A3) -- (A2); \draw[hid] (A3) -- (B3);
    % The cut runs BACK corner -> FRONT corner of the top face (B3 -> B1).  The
    % old B0--B2 diagonal ran almost along the top edge, so it read as a stray
    % duplicate edge; and the book draws no companion diagonal on the base.
    \draw[fig] (B3) -- (B1);
    \node[slab, below] at (1.72,-0.55) {(a)};
  \end{scope}
  % (b) an arbitrary prism, cut into triangular prisms
  \begin{scope}[xshift=5.35cm]
    \coordinate (ea) at (2.95,0); \coordinate (eb) at (1.02,0.72);
    \coordinate (ec) at (0,1.42);
    \coordinate (A0) at (0,0);   \coordinate (A1) at ($(A0)+(ea)$);
    \coordinate (A2) at ($(A1)+(eb)$); \coordinate (A3) at ($(A0)+(eb)$);
    \coordinate (B0) at ($(A0)+(ec)$); \coordinate (B1) at ($(A1)+(ec)$);
    \coordinate (B2) at ($(A2)+(ec)$); \coordinate (B3) at ($(A3)+(ec)$);
    \draw[fig] (A0) -- (A1) -- (A2); \draw[fig] (A0) -- (B0);
    \draw[fig] (A1) -- (B1); \draw[fig] (A2) -- (B2);
    \draw[fig] (B0) -- (B1) -- (B2) -- (B3) -- cycle;
    \draw[hid] (A0) -- (A3) -- (A2); \draw[hid] (A3) -- (B3);
    % a vertical cutting plane
    \coordinate (S0) at ($(A0)!0.62!(A1)$);
    \coordinate (S1) at ($(B0)!0.62!(B1)$);
    \coordinate (S2) at ($(B3)!0.62!(B2)$);
    \draw[fig] (S0) -- (S1) -- (S2);
    % the cut's hidden edge is the VERTICAL down the back face, not a line
    % across the base: it drops from the top-back edge to the hidden bottom one
    \draw[hid] (S2) -- ($(A3)!0.62!(A2)$);
    \node[slab, below] at (2.02,-0.55) {(b)};
  \end{scope}
\end{tikzpicture}
\figcap{14--38}
\end{bkfigure}}

% --------------------------------------------------------------- Fig 14-39
\newcommand{\FIGXIVTHIRTYNINE}{\begin{bkfigure}[0.50\linewidth]
\begin{tikzpicture}[scale=1.560]
  % base is a parallelogram; ct is where its diagonals meet, so the altitude
  % really does fall on the centre of the base.
  % The base is a parallelogram whose front vertex is NOT under the apex, so
  % the altitude reads as its own line instead of hiding inside the front
  % edge -- which is how the book draws it.
  \coordinate (b1) at (0,0.55); \coordinate (b2) at (1.30,0);
  \coordinate (b3) at (3.05,0.55); \coordinate (b4) at (1.75,1.10);
  \coordinate (ct) at (1.525,0.55);          % centre = crossing of the diagonals
  \coordinate (ap) at (1.525,2.85);
  \draw[fig] (b1) -- (b2) -- (b3);
  \draw[hid] (b3) -- (b4) -- (b1);
  \draw[fig] (ap) -- (b1); \draw[fig] (ap) -- (b2); \draw[fig] (ap) -- (b3);
  \draw[hid] (ap) -- (b4);
  \draw[aux, ->, >=latex] (ap) -- (ct);
  \node[vlab] at (1.74,0.92) {$h$};
\end{tikzpicture}
\figcap{14--39}
\end{bkfigure}}

% --------------------------------------------------------------- Fig 14-40
% Two similar pyramids: the right one is the left one scaled by 1.55 about V'.
\newcommand{\FIGXIVFORTY}{\begin{bkfigure}[0.66\linewidth]
\begin{tikzpicture}[scale=0.780]
  % ---- small pyramid, vertex V at the lower left of the base
  \begin{scope}
    \coordinate (V)  at (0,0);
    \coordinate (v2) at (1.35,0);
    \coordinate (v3) at (1.92,0.42);
    \coordinate (ap) at (0.78,1.62);
    \draw[fig] (V) -- (v2) -- (v3);
    \draw[fig] (ap) -- (V); \draw[fig] (ap) -- (v2); \draw[fig] (ap) -- (v3);
    \node[vlab, below=3.5pt] at (V) {$V$};
    \node[vlab, above right=3pt] at ($(ap)!0.45!(v3)$) {$l$};
  \end{scope}
  % ---- large pyramid, the same shape scaled by 1.55
  \begin{scope}[xshift=2.75cm]
    \coordinate (Vp)  at (0,0);
    \coordinate (w2) at (2.0925,0);
    \coordinate (w3) at (2.976,0.651);
    \coordinate (bp) at (1.209,2.511);
    \draw[fig] (Vp) -- (w2) -- (w3);
    \draw[fig] (bp) -- (Vp); \draw[fig] (bp) -- (w2); \draw[fig] (bp) -- (w3);
    \node[vlab, below=3.5pt] at (Vp) {$V'$};
    \node[vlab, above right=3pt] at ($(bp)!0.45!(w3)$) {$l'$};
  \end{scope}
\end{tikzpicture}
\figcap{14--40}
\end{bkfigure}}

% --------------------------------------------------------------- Fig 14-41
\newcommand{\FIGXIVFORTYONE}{\begin{bkfigure}[0.82\linewidth]
\begin{tikzpicture}[scale=0.858]
  % ---- cylinder, tilted
  \begin{scope}[rotate=-18]
    \draw[fig] (0,0) ellipse (0.62 and 0.24);
    \draw[fig] (-0.62,0) -- (-0.62,-2.35);
    \draw[fig] (0.62,0) -- (0.62,-2.35);
    \draw[fig] (-0.62,-2.35) arc (180:360:0.62 and 0.24);
    \draw[hid] (0.62,-2.35) arc (0:180:0.62 and 0.24);
  \end{scope}
  % ---- cone
  \begin{scope}[xshift=2.85cm, yshift=-2.05cm]
    \draw[fig] (-0.78,0) arc (180:360:0.78 and 0.28);
    \draw[hid] (0.78,0) arc (0:180:0.78 and 0.28);
    \draw[fig] (-0.78,0) -- (0,2.05) -- (0.78,0);
  \end{scope}
  % ---- sphere, with the equator
  \begin{scope}[xshift=5.75cm, yshift=-1.02cm]
    \draw[fig] (0,0) circle (1.08);
    \draw[fig] (-1.08,0) arc (180:360:1.08 and 0.38);
    \draw[hid] (1.08,0) arc (0:180:1.08 and 0.38);
  \end{scope}
\end{tikzpicture}
\figcap{14--41}
\end{bkfigure}}

% ---------------------------------------------------------------- Fig 14-42
% Split out of the old \FIGXIVFORTYTWOFORTYTHREE block: Ex.~14-16 \#2 cites
% 14-42, while 14-43 belongs to the body text about Theorem 14-24.
\newcommand{\FIGXIVFORTYTWO}{\begin{bkfigure}[0.50\linewidth]
\begin{tikzpicture}[scale=1.35]
  % Right circular cone with an inscribed pyramid.  The book's base polygon is
  % a HEXAGON with vertices at the two extremes of the ellipse plus two in
  % front and two behind: the three near base chords are solid, the three far
  % ones dashed, and the two hidden lateral edges stand clear of the solid
  % ones instead of smearing into them at the apex.  (Ours drew a single
  % dashed chord and no near chords at all, so nothing read as a pyramid, and
  % it carried a spurious dashed diameter across the base.)
  \coordinate (ap) at (0,2.07);
  \draw[fig] (0,0) ellipse (1.42 and 0.61);
  \foreach \i/\a in {0/180, 1/230, 2/310, 3/0, 4/50, 5/130}
    {\coordinate (e\i) at ({1.42*cos(\a)},{0.61*sin(\a)});}
  \draw[fig] (e0) -- (e1) -- (e2) -- (e3);          % near base chords
  \draw[hid] (e3) -- (e4) -- (e5) -- (e0);          % far base chords
  \foreach \i in {0,1,2,3}{\draw[fig] (ap) -- (e\i);}
  \foreach \i in {4,5}{\draw[hid] (ap) -- (e\i);}
\end{tikzpicture}
\figcap{14--42}
\end{bkfigure}}

% ---------------------------------------------------------------- Fig 14-43
\newcommand{\FIGXIVFORTYTHREE}{\begin{bkfigure}[0.48\linewidth]
\begin{tikzpicture}[scale=0.980]
  % a sphere carved into small "pyramids" with a common apex at the centre
  \begin{scope}[even odd rule]
    \clip (0,0) circle (1.42);
    \clip (0.28,-0.24) circle (1.28) (-1.6,-1.6) rectangle (1.6,1.6);
    \fill[pattern=north east lines] (-1.6,-1.6) rectangle (1.6,1.6);
  \end{scope}
  \draw[fig] (0,0) circle (1.42);
  % one small pyramid, apex at the centre
  \coordinate (O) at (0,0);
  \coordinate (k1) at (0.72,0.42); \coordinate (k2) at (1.02,0.62);
  \coordinate (k3) at (1.02,0.05); \coordinate (k4) at (0.72,-0.10);
  \draw[fig] (k1) -- (k2) -- (k3) -- (k4) -- cycle;
  \draw[fig] (O) -- (k1); \draw[fig] (O) -- (k4);
  \draw[hid] (O) -- (k2); \draw[hid] (O) -- (k3);
  \dt{O}
\end{tikzpicture}
\figcap{14--43}
\end{bkfigure}}
